% Template for post or presentation abstracts submitted to ILAS 2022
% 1. Fill in the presenter's information (name, affiliation, email
% address) and talk information (title of presentation, abstract).
% 2. Compile to make sure there are no errors.
% 3. Save the .tex file for your abstract with a distinctive name,
% ideally "FamilynameGivenname.tex".
% 4. Upload your .tex file to https://clr.ie/129557
%       
% * Please keep your LaTeX commands simple, and avoid the use of
%    user-defined macros.
% * Include details of co-authors and funding acknowledgements after the abstract.
% * Upload only the LaTeX file (with .tex extension).
% * Questions: Contact Niall Madden (Niall.Madden@NUIGalway.ie)

\documentclass[a4paper]{amsart} 
\usepackage{amssymb} 
\usepackage{hyperref}
\pagestyle{empty}
\date{} 

%% IGNORE THE NEXT 4 LINES
\makeatletter % 
\def\labelenumi{\theenumi.}
\def\theenumi{\@arabic\c@enumi}
\makeatother
%% STOP IGNORING NOW

\begin{document}

\title{Eigenvalues, spanning trees, and connected parity factors in regular graphs}
\author{Suil O} %Presenting author only
\address{SUNY-Korea} % Affiliation of presenting author
\email{suil.o@sunykorea.ac.kr} % Email address of presenting author only

\maketitle

% The abstract  ***no more than one page***

In this talk, we prove a sharp upper bound for the second largest eigenvalue in an $r$-regular graph $G$ to guarantee that $G$ contains at least two disjoint spanning trees. By utilizing the result, we prove an upper bound for the second largest eigenvalue in an $r$-regular graph to guarantee the existence of a connected parity factor.

%% Coauthor details here (optional - delete if not applicable)
% and funding acknowledgment (optional - delete if not applicable)
\emph{This is joint work with Donggyu Kim (KAIST) and Zhiwen Wang (Nankai).
Supported by the National Research Foundation of Korea, Grant NRF-2020R1F1A1A01048226.}

\begin{thebibliography}{1}
\bibitem{KO}
Donggyu Kim and Suil O.
\newblock Eigenvalues and parity factors in graphs.
\newblock {\em  Submitted}.

\bibitem{KOPR}
Sungeun Kim, Suil O, Jihwan Park, and Hyo Ree.
\newblock An odd $[1,b]$-factor in regular graphs from eigenvalues.
\newblock {\em  Disc. Math.}  343:111906,  (2020).

\bibitem{KOW}
Donggyu Kim, Suil O, and Zhiwen Wang.
\newblock Eigenvalues, spanning trees, and connected parity factors in regular graphs.
\newblock {\em  (in preparation)}.

\bibitem{O}
Suil O.
\newblock Eigenvalues and [a,b]-factors in regular graphs.
\newblock {\em  J. Graph Theory} (in press).
\end{thebibliography}


\end{document}


% Please leave the next bit alone  
%%% Local Variables: 
%%% mode: latex
%%% TeX-master: "../ILAS-2022"
%%% End:
